\documentclass[12pt, a4paper]{article}

\usepackage{fontspec}
\usepackage{polyglossia}
\setdefaultlanguage{bengali}
\setotherlanguage{english}

% Provide a Bengali font. 'Kalpurush' or 'SolaimanLipi' are commonly used.
% You must compile this file with XeLaTeX or LuaLaTeX. 
\newfontfamily\bengalifont[Script=Bengali, AutoFakeBold=2.5, AutoFakeSlant=0.2]{Kalpurush}

\usepackage[margin=1in]{geometry}
\usepackage{hyperref}
\usepackage{color}
\usepackage{xcolor}
\usepackage{tcolorbox}
\usepackage{listings}
\usepackage{enumitem}

\hypersetup{colorlinks=true, linkcolor=blue, urlcolor=cyan, pdftitle={Native Core Analysis}}

\definecolor{codegreen}{rgb}{0,0.6,0}
\definecolor{codegray}{rgb}{0.5,0.5,0.5}
\definecolor{codepurple}{rgb}{0.58,0,0.82}
\definecolor{backcolour}{rgb}{0.95,0.95,0.92}

\lstdefinestyle{mystyle}{
    backgroundcolor=\color{backcolour},   
    commentstyle=\color{codegreen},
    keywordstyle=\color{magenta},
    numberstyle=\tiny\color{codegray},
    stringstyle=\color{codepurple},
    basicstyle=\ttfamily\small,
    breakatwhitespace=false,         
    breaklines=true,                 
    captionpos=b,                    
    keepspaces=true,                 
    numbers=left,                    
    numbersep=5pt,                  
    showspaces=false,                
    showstringspaces=false,
    showtabs=false,                  
    tabsize=4
}
\lstset{style=mystyle}

\title{\textbf{\Large Native Core Analysis\\ (Constants, AccessibilityUtils \& OverlayManager)}}
\author{Synapse / No To Distraction}
\date{\today}

\begin{document}

\maketitle

\section*{ভূমিকা}
আপনার দেওয়া ৩টি ফাইল (\texttt{Constants.kt}, \texttt{AccessibilityUtils.kt}, এবং \texttt{OverlayManager.kt}) মূলত অ্যাপের কোর লজিক, সেটিংস এবং ব্লকিং স্ক্রিন দেখানোর কাজে ব্যবহৃত হয়। প্রতিটি ফাইলের ইমপোর্ট, ফাংশন, এবং এবং অন্যান্য ফাইলের সাথে এদের কানেকশন বিস্তারিতভাবে নিচে দেওয়া হলো:

\vspace{0.5cm}

\begin{tcolorbox}[colback=blue!5!white,colframe=blue!75!black,title={১. Constants.kt}]
এই ফাইলটিতে অ্যাপের সমস্ত \textbf{ধ্রুবক (Constants)} বা ফিক্সড ভ্যালুগুলো স্টোর করে রাখা হয়েছে। যাতে পুরো অ্যাপে বারবার একই নাম বা নম্বর টাইপ করতে না হয়। এটি মূলত একটি ডিকশনারির মতো কাজ করে।

\textbf{ইমপোর্টস (Imports):}
\begin{itemize}[label=$\circ$]
    \item \texttt{AccessibilityEvent}: অ্যান্ড্রয়েডের স্ক্রিন ইভেন্টগুলো চেনার জন্য।
\end{itemize}

\textbf{মূল ব্লক ও ভেরিয়েবলসমূহ:}
\begin{itemize}[label=$\circ$]
    \item \textbf{SharedPreferences Keys:} \texttt{PREF\_BLOCK\_FB\_REELS}, \texttt{PREF\_FOCUS\_MODE\_END\_TIME\_MS} ইত্যাদি। এগুলো হলো মেমোরিতে ডেটা সেভ করার "চাবি" (Key)। 
    \begin{itemize}[label=--]
        \item \textit{কানেকশন:} \texttt{MainActivity.kt} এবং \texttt{StorageManager.kt} ফাইল এই চাবিগুলো ধরে ডেটা সেভ করে ও পড়ে।
    \end{itemize}
    \item \textbf{Package Names:} \texttt{FACEBOOK\_PACKAGE}, \texttt{INSTAGRAM\_PACKAGE}: ফেসবুক, ইন্সটাগ্রাম ও ইউটিউবের আসল প্যাকেজ নেইম এখানে ফিক্সড করা।
    \item \textbf{Thresholds \& Durations:} 
    \begin{itemize}[label=--]
        \item \texttt{EVENT\_DEBOUNCE\_MS (400L)}: খুব দ্রুত বারবার ইভেন্ট আসা ঠেকাতে ৪০০ মিলিসেকেন্ড ডিলে।
        \item \texttt{REELS\_LOCK\_DURATION\_MS}: রিল ব্লক ফিচারটি একবার লক করলে কতক্ষণ লক থাকবে তার হিসেব।
    \end{itemize}
    \item \textbf{Spatial Heuristic Constants:} এটি খুবই ইন্টারেস্টিং! রিল বা শর্টস ভিডিওগুলো সাধারণত পুরো স্ক্রিন জুড়ে থাকে। এই রেশিওগুলো দিয়ে \texttt{ScannerEngine.kt} গাণিতিক পরিমাপে বুঝতে পারে ইউজারের স্ক্রিনে চলা ভিডিওটি আসলেই 'রিলস' কি না। 
    \item \textbf{TARGET\_EVENT\_TYPES:} সার্ভিসটি শুধুমাত্র স্ক্রিনের উইন্ডো চেঞ্জ বা স্ক্রোলিং ইভেন্টগুলোকে মনিটর করবে তার একটি লিস্ট।
\end{itemize}
\end{tcolorbox}

\begin{tcolorbox}[colback=green!5!white,colframe=green!75!black,title={২. AccessibilityUtils.kt}]
অ্যাক্সেসিবিলিটি সার্ভিসটি (যা স্ক্রিন স্ক্যান করে) ঠিকমতো চালু আছে কি না, তা একদম নির্ভুলভাবে চেক করবে এই ফাইল।

\textbf{ইমপোর্টস (Imports):}
\begin{itemize}[label=$\circ$]
    \item \texttt{Context, ComponentName, Settings, AccessibilityManager, ContextCompat}: সিস্টেমের সার্ভিস লিস্ট এবং সেটিংস চেক করার জন্য প্রয়োজনীয় অ্যান্ড্রয়েড টুলস।
\end{itemize}

\textbf{ফাংশন এবং কানেকশন:}
\begin{itemize}[label=$\circ$]
    \item \textbf{\texttt{isShortVideoServiceEnabled()}}
    \begin{itemize}[label=--]
        \item \textbf{কাজ:} ফাংশনটি দুই ধাপে চেকিং করে। প্রথমে এটি অ্যান্ড্রয়েডের \texttt{AccessibilityManager} কে জিজ্ঞেস করে। অনেক সময় ফোনে ক্যাশ এর কারণে এটি ভুল উত্তর দেয়, তাই দ্বিতীয় ধাপে এটি সরাসরি ফোনের সেটিংস এর কোর ডেটাবেস (\texttt{Settings.Secure}) চেক করে কনফার্ম হয়। 
        \item \textbf{নেটিভ কানেকশন:} \texttt{MainActivity.kt} এবং \texttt{ShortVideoAccessibilityService.kt} এখান থেকে এটি কল করে।
        \item \textbf{Flutter কানেকশন:} ফ্লাটারের \texttt{accessibility\_permission\_service.dart} এই স্ট্যাটাসটিই দেখতে পায়।
    \end{itemize}
\end{itemize}
\end{tcolorbox}

\begin{tcolorbox}[colback=orange!5!white,colframe=orange!75!black,title={৩. OverlayManager.kt}]
অ্যাপের সবচাইতে ভিজ্যুয়াল বা দৃশ্যমান কাজ করে এই ফাইলটি। যখনই কোনো রিলস বা ব্লক করা অ্যাপ ওপেন হয়, তখন উইন্ডোটি বা ব্লকার ভেসে ওঠে, তা এই ফাইলটিই তৈরি করে।

\textbf{ইমপোর্টস (Imports):}
\begin{itemize}[label=$\circ$]
    \item \texttt{WindowManager, View, Button, TextView} ইত্যাদি: উইন্ডো, ওভারলে ড্র করার জন্য।
    \item \texttt{Handler, Looper, SystemClock}: লাইভ টাইমার বা কাউন্টডাউন চালানোর জন্য।
\end{itemize}

\textbf{ফাংশন ও ব্লকের কাজ:}
\begin{itemize}[label=$\circ$]
    \item \textbf{ক্লাস লেভেল ভেরিয়েবল:} \texttt{windowManager} হলো অ্যান্ড্রয়েডের এমন এক টুল যা যেকোনো অ্যাপের উপরে উইন্ডো ভাসিয়ে রাখতে পারে। \texttt{countdownRunnable} হলো একটি লুপ যা প্রতি সেকেন্ডে একবার করে কল হয়ে টাইমার আপডেট করে।
    
    \item \textbf{\texttt{showReelsOverlay(isFocusMode: Boolean)}}
    \begin{itemize}[label=--]
        \item \textbf{কাজ:} যখন ইউজার রিলস দেখে, তখন এটি পুরো স্ক্রিনের উপর একটি কালো লেয়ার বসিয়ে দেয়। ফোকাস মোড চলাকালীন সময়ে এটি একটি সুন্দর মোটিভেশনাল উক্তিও দেখায়।
        \item \textbf{কানেকশন:} \texttt{ScannerEngine.kt} যখন রিলস ডিটেক্ট করে তখন এটি কল হয়।
        \item \textbf{Flutter কানেকশন:} এই ফাংশনের ভেতর থেকেই \texttt{ReelDetectionChannelBridge} কে কল করা হয়, যা সাথে সাথেই ফ্লাটার অ্যাপকে জানিয়ে দেয় একটি রিল ব্লক করা হয়েছে।
    \end{itemize}

    \item \textbf{\texttt{showQuickBlockOverlay(...)}} ও \textbf{\texttt{hideQuickBlockOverlay()}}
    \begin{itemize}[label=--]
        \item \textbf{কাজ:} এটি আলাদা একটি ওভারলে। যখন পুরো একটি অ্যাপ কুইক ব্লকের মাধ্যমে ব্লক করা থাকে, তখন এই ব্লক স্ক্রিনটি দেখায় এবং একটি লাইভ টাইমার চলতে থাকে।
    \end{itemize}

    \item \textbf{\texttt{isSuppressed()}}
    \begin{itemize}[label=--]
        \item \textbf{কাজ:} ইউজার "Go Home" এ ক্লিক করলে এই ফাংশনটি আড়াই সেকেন্ডের জন্য ব্লকারকে মিউট রাখে যাতে ওভারলেটি হোমস্ক্রিনে বারবার ফ্লিকার বা গ্লিচ না করে। 
    \end{itemize}
    
    \item \textbf{\texttt{createLayoutParams()}}
    \begin{itemize}[label=--]
        \item \textbf{কাজ:} উইন্ডোটি ఎంతটুক বড় হবে, তা কনফিগার করে। এখানে \texttt{TYPE\_APPLICATION\_OVERLAY} ব্যবহার করা হয় যাতে এটি সিস্টেমের একদম সবকিছুর উপরে থাকতে পারে।
    \end{itemize}
\end{itemize}
\end{tcolorbox}

\vspace{0.5cm}
\textbf{সংক্ষেপে (Summary Flow):}\\
\texttt{ScannerEngine.kt} (রিলস খুঁজলো) $\rightarrow$ রুলস মেলালো \texttt{Constants.kt} থেকে $\rightarrow$ রিলস পেলে \texttt{OverlayManager.kt} কে নির্দেশ দিলো ব্লকার দেখাতে $\rightarrow$ ওভারলে স্ক্রিনে ভেসে উঠলো $\rightarrow$ ফ্লাটারকে জানিয়ে দেওয়া হলো।

\end{document}
